\documentclass[10pt,a4paper]{amsart}
\usepackage[margin=19mm]{geometry}
\usepackage{amsmath,amssymb,mathtools}
\usepackage[scr=rsfs,cal=euler]{mathalfa}
\usepackage{xfrac}
\usepackage[unicode,hidelinks,bookmarks=false]{hyperref}
\newcommand{\ie}{i.e.\ }
\newcommand{\cf}{cf.\ }
\newcommand{\resp}{respectively\ }
\newcommand{\Subsection}{\subsection}
\providecommand{\nobreakdash}{\nobreak\hbox{-}}
\setlength{\emergencystretch}{2em}
\multlinegap0pt
\usepackage[normalem]{ulem}
\usepackage{amssymb}
\usepackage[shortlabels]{enumitem}
\newtheorem{theorem}{Theorem}
\newtheorem{lem}{Lemma}
\def \F{{\mathbb F}}
\def \P{{\mathcal P}}
\def \al{\alpha}
\def \avsp{\mbox{\bf avsp}}
\def \be{\beta}
\def \de{\delta}
\def \ga{\gamma}
\def \la{\left\langle}
\def \o{\mbox{\bf 0} }
\def \ra{\right\rangle}
\title{A note on Tight Irreducible Affine Spreads}
\author{Fusun Akman, Papa Sissokho}
\date{Source: arXiv:2607.23411v1 (CC BY 4.0)}
\begin{document}
\maketitle

\noindent\emph{Source and licence.} A note on Tight Irreducible Affine Spreads. Version arXiv:2607.23411v1 (CC BY 4.0), distributed by arXiv under the Creative Commons Attribution 4.0 International licence, http://creativecommons.org/licenses/by/4.0/, as recorded in the arXiv licence metadata for this exact version and shown on the abstract page https://arxiv.org/abs/2607.23411v1. Full licence notice: \texttt{licenses/arxiv-2607.23411-CC-BY-4.0.txt}.\par\smallskip

\noindent\textit{Editorial scope.} Counted unit: the article's theorem thm:2 (source0.tex lines 187--189). The article's complete original proof of it is delivered with it (source0.tex lines 191--237, one \texttt{proof} environment of 47 lines). No mathematics is written, repaired or re-derived: the statement and the proof are the article's own lines, in the article's own notation, and no retained line is substituted or re-flowed.  Retained uncounted as the source-local context the proof needs: lines 104--106; lines 139--141.  Nothing was omitted from any retained span, so the package carries no omission record.  No retained line contains a citation, so the wrapper carries no bibliography and no bibliography entry.  No retained line contains a figure environment, an image-inclusion command or drawing code, so no image is delivered and the package declares no asset dependency.  Editorial formatting only, in addition: an AMS \texttt{amsart} preamble that restates the article's own declarations and macros used by the retained lines, namely the environments \texttt{theorem}, \texttt{lem} on separate counters, together with the standard packages those lines need.  The retained environments print in this document as Theorem 1, Lemma 1 in order of appearance, whereas the article prints the same items with the numbering of its own full document.  The complete native source of the article as distributed in the arXiv e-print for this exact version is bundled unchanged as \texttt{sources/206025/source0.tex}.
\begin{theorem}\label{thm:1}
If $d \geq 1$ and $n>2d$, then there exists a completely tight irreducible affine $d$-spread of AG$(n,q)$.  
\end{theorem}

\begin{lem}\label{lem:1}
The set $\P_\al$  given in Construction~1 is  a completely tight affine $d$-spread of  AG$(n,q)$.
\end{lem}

\begin{theorem}\label{thm:2}
Let $\P$ be the affine partition given in Construction~2. Then $\P$ is a completely tight irreducible affine $d$-spread of AG$(n,q)$.
\end{theorem}

\begin{proof}{Proof of Theorem \ref{thm:1}}
Since Construction~2 retains all the relevant properties of Construction~1, the fact that $\P=\{C_{\ga}:\; \ga\in H\}$ is a completely tight \avsp~follows from Lemma~\ref{lem:1}.

Next, we prove the irreducibility of $\P$. Assume that there exists some sub-partition $\P'\subset \P$ such that $1<|\P'|<|\P|$ and the union of all subspaces in $\P'$ is a subspace $V'$ of AG$(n,q)$. 
By translating\footnote{(Complete) tightness and irreducibility of $\P_\al$'s are preserved by vector translations.} $\P$ if necessary, we can assume that $V'$ is linear. Thus, if $\P$ is reducible, then we have a proper subspace $V'$ of $V$ such that
%
\begin{equation}\label{eq:span}
V'=\la \bigcup_{C\in \P'} C\ra=\bigcup_{C\in \P'} C.
\end{equation}
%
Since $\P'$ contains at least two cosets and $V'$ is a linear subspace,  one of the minimum of two cosets must be $C_{0}=S_0=U$ and the other one $C_{\ga}$, for some $\ga\not=0$.
Then it follows from~\eqref{eq:span} that
\[U=\la u_1, \ldots, u_d\ra \subseteq V'\mbox{ and } \la C_{\ga}\ra = \la \ga, \ga\be+u_1,\ldots,\ga\be^d+u_d\ra\subseteq V',\]
and hence,
 \[\la \ga,\ga\be,\ldots,\ga\be^d, u_1, \ldots, u_d\ra \subseteq V'.\]
 We now show by induction that $\ga \be^k\in V'$ for all integers $k\geq1$.
Given that $\ga \be\in V'$, the base case holds. Next, assume that $\ga \be^k\in V'$ for any $k\geq1$. Since $\ga\in V'$ and $V'$ is a linear subspace, we must have $\ga+\ga \be^k\in V'$. Thus, 
it follows from~\eqref{eq:span} that there exists $\de\in H$ such that  $C_\de\in\P'$ and $\ga+\ga \be^k \in C_\de$. Consequently, we have
%
\begin{align*}\label{eq:cond1}
\ga+\ga \be^k \in C_\de 
& \Longleftrightarrow 
\ga+\ga \be^k = \de+\sum_{i=1}^d a_i (\de\be^i+u_i) \quad \mbox{(for some scalars $a_i\in \F_q$)}\cr 
& \Longleftrightarrow 
-\de+\ga+\ga \be^k -\sum_{i=1}^d a_i \de\be^i=\sum_{i=1}^d a_i u_i  \cr 
& \Longrightarrow 
 \mbox{ $a_i=0$ for $1\leq i\leq d$ } \quad (H\cap U=\{ \o\}) 
 \cr
 & \Longrightarrow 
\de=\ga+\ga \be^k.
\end{align*}
%

 \[C_\de= C_{\ga+\ga \be^k}=\ga+\ga \be^k+ \la (\ga+\ga \be^k)\be+u_1,\ldots,(\ga+\ga \be^k)\be^d+u_d\ra \subseteq V'.\]
In particular,  it must be true that
%
\begin{align}\label{eq:cond2}
\ga+\ga \be^k+\left((\ga+\ga \be^k)\be+u_1\right) =
\ga+\ga\be+\ga \be^k+\ga \be^{k+1}+u_1 \mbox{ is in } V'\Longrightarrow \ga \be^{k+1}\in V' .
\end{align}
%
This shows the inductive step and proves that $\ga \be^k\in V' $ for any integer $k\geq 1$.
Since $\ga\not=0$ and $H=\F_q(\beta) \cong \F_{q^{n-d}}$, it follows that 
\[\{\ga \be^k:\; k\geq 1\}=\ga\F_{q^{n-d}}^*=\F_{q^{n-d}}^*\subseteq V',\]
where $\F_{q^{n-d}}^*=\F_{q^{n-d}}\setminus\{\o\}$. 
Moreover $U\subseteq V'$, and thus, $V=H\oplus U\subseteq V'$, which implies that $V'=V$. We have proven the irreducibility of $\P=\{C_\ga:\; \ga\in H\}$ and the proof is complete.
\end{proof} 

\end{document}
